% Obsahový zdroj: import_hodiny/35_Kombinace_bez_opakovani.tex a odpovídající PDF (obě strany).
% P1–P5: zdrojové Příklady 1–5; C1–C5: zdrojové Procvičování 1–5.
% Matematický obsah zachován; výklad, odvození a řešení pedagogicky rozvedeny.
% Upřesnění a opravy: viz KMB_29-41_Kontrola.md; původní import zůstává beze změny.
% Rozpis: 4_rocnik/tematicke_plany/Rozpis_vyuky_MAT_4_rocnik.tex, hodina 35.
% Plán: 4_rocnik/tematicke_plany/SPS.pdf, oddíly 2.1–2.3.
% Sazba: sablony/pripravy/TEMPLATE_Priprava.tex a společný priprava.sty.
% Vlastní doplnění: didaktický výklad, zdůvodnění, metodika; žádné nové číslované úlohy.
% Učebnice/09_Kombinatorika nemá obsahovou kapitolu; místní kombinatorické listy ani Realisticky.cz nejsou dostupné.
\documentclass[11pt,a4paper]{article}
% Spolecny styl pro ucitelske pripravy MAT (Work/MAT)
% Verze 1.1, 2026-10-23
\NeedsTeXFormat{LaTeX2e}
\ProvidesPackage{priprava}[2026/10/23 v1.1 Sablona ucitelskych priprav MAT]

% Zakladni typografie a cestina
\RequirePackage[T1]{fontenc}
\RequirePackage[utf8]{inputenc}
\RequirePackage[czech]{babel}
\RequirePackage{lmodern}
\RequirePackage{microtype}
\RequirePackage{amsmath,amssymb}
\RequirePackage{xcolor}
\RequirePackage{enumitem}
\RequirePackage[most]{tcolorbox}
\RequirePackage{tabularx,booktabs,array}
\RequirePackage{fancyhdr}
\RequirePackage{needspace}
\RequirePackage{geometry}
\RequirePackage[hidelinks]{hyperref}

\geometry{a4paper,top=17mm,bottom=18mm,left=23mm,right=23mm}
\definecolor{notegray}{gray}{0.39}
\setlength{\parindent}{0pt}
\setlength{\parskip}{2.5pt plus 1pt minus 0.5pt}
\setlist[itemize]{topsep=2pt,itemsep=1pt,parsep=0pt,leftmargin=1.45em}
\setlist[enumerate]{topsep=2pt,itemsep=2pt,parsep=0pt,leftmargin=1.65em}
\setcounter{secnumdepth}{0}
\renewcommand{\arraystretch}{1.12}

% Zahlavi: vlevo rocnik | globalni tema, vpravo cislo: nazev hodiny
% Zapati zustava zcela prazdne.
\pagestyle{fancy}
\fancyhf{}
\renewcommand{\headrulewidth}{0.35pt}
\renewcommand{\footrulewidth}{0pt}
\setlength{\headheight}{14pt}
\newcommand{\nastavPripravu}[4]{%
  \fancyhead[L]{\small #1\enspace\textbar\enspace #2}%
  \fancyhead[R]{\small #3: #4}%
}

% Minimalisticke cernobile boxy. Kazda definice zvlast.
\tcbset{commonlesson/.style={enhanced,breakable=false,colback=white,colframe=black,
  colbacktitle=white,coltitle=black,boxrule=0.45pt,
  arc=0.8mm,outer arc=0.8mm,left=3mm,right=3mm,top=1.6mm,bottom=1.6mm,
  toptitle=1mm,bottomtitle=0.7mm,fonttitle=\bfseries,
  attach boxed title to top left={xshift=2.5mm,yshift=-1.15mm},
  boxed title style={colback=white,colframe=white,boxrule=0pt,left=1mm,right=1mm}}}
\newtcolorbox{zadaniBox}[1][]{commonlesson,title={Zadání},#1}
\newtcolorbox{prikladBox}[1][]{commonlesson,title={Příklad},#1}
\newtcolorbox{definiceBox}[1][]{commonlesson,title={Definice},#1}

% Spolecny box pro zadani a jeho reseni, bez mezery mezi castmi.
% Pouziti:
%   \begin{prikladReseniBox}
%     Zadani ulohy.
%     \tcblower
%     \begin{align*} ... \end{align*}
%   \end{prikladReseniBox}
% Spodni cast ma jemny sedy podklad a delici linku; lze v ni sazet
% vice radku matematiky. Obe casti maji jen jeden spolecny ramecek.
\tcbset{solutionlesson/.style={
  commonlesson,
  enhanced jigsaw,
  breakable,
  colbacklower=black!3,
  segmentation style={solid,draw=black!25,line width=0.35pt},
  fontlower=\small\color{notegray},
  before lower={\noindent\textbf{Řešení:}\enspace},
  leftlower=3mm,rightlower=3mm
}}
\newtcolorbox{prikladReseniBox}[1][]{solutionlesson,title={Příklad},#1}
\newtcolorbox{zadaniReseniBox}[1][]{solutionlesson,title={Zadání},#1}

% Pomucky pro autora pripravy: sedive poznamky; \reseni pro starsi pripravy
\newcommand{\poznamka}[1]{\par\begingroup\small\color{notegray}\itshape
  \textnormal{\textbf{Poznámka pro učitele.}}\enspace #1\par\endgroup}
\newcommand{\reseni}[1]{\par\begingroup\small\color{notegray}
  \textbf{Očekávané řešení:}\enspace #1\par\endgroup}
\newcommand{\blok}[1]{\needspace{5\baselineskip}\subsection*{#1}}
\newcommand{\vyrok}[1]{\uv{#1}}

\nastavPripravu{4. ročník}{Kombinatorika}{35}{Kombinace bez opakování}
\begin{document}


Připomeň obsazení tří různých funkcí. Nyní místo funkcí vybíráme pouze členy komise: stejní lidé mají tvořit jediný výsledek bez ohledu na pořadí, v němž byli vybráni.
\poznamka{Předpoklady: variace bez opakování, faktoriál, součin voleb. Cíl: žák vysvětlí dělení $k!$ a použije kombinace ve slovních a geometrických úlohách. Začni P1, pak P2 a P3. P4–P5 jsou rozšíření aplikační části; při pomalejším tempu společně zvol jen jeden z nich, druhý zůstává rezervou.}

\blok{Výklad a odvození}
\begin{definiceBox}
\textbf{Kombinace bez opakování} je neuspořádaný výběr $k$ různých prvků z $n$ navzájem různých prvků. Pro celá $n\geq0$, $0\leq k\leq n$ je jejich počet $C(n,k)$; zapisujeme jej také $\binom nk$ a čteme „$n$ nad $k$“.
\end{definiceBox}

U každé konkrétní skupiny $k$ lidí lze pořadí jejich výběru zapsat $k!$ způsoby. Variace $V(n,k)$ tedy každou skupinu počítají právě $k!$-krát. Například jedna trojice A, B, C dává pořadí
\[
 ABC,\ ACB,\ BAC,\ BCA,\ CAB,\ CBA.
\]
Aby se trojice započetla jednou, dělíme šesti, nikoli třemi. Obecně
\[
 C(n,k)=\frac{V(n,k)}{k!}=\frac{n!}{k!(n-k)!}=\binom nk.
\]
Odvození je platné, protože \emph{každá} skupina má tentýž počet uspořádání. Pro $k=0$ i $k=n$ je jeden výběr: prázdná množina, resp. celá množina. Při výběru s opakováním se tento vzorec nepoužije; takové výběry dnes neřešíme.

Rozhodovací otázka zní: „Změní záměna dvou vybraných lidí výsledek?“ V komisi bez funkcí nezmění; při obsazení předsedy a místopředsedy změní. U geometrických úloh navíc prověřujeme, zda každá vybraná skupina bodů skutečně určí požadovaný objekt. Dva různé body určují úsečku, tři kolineární body však neurčují trojúhelník.

Má-li skupina obsahovat předepsané počty členů z několika disjunktních kategorií, provedeme jednotlivé výběry a jejich počty násobíme. Celá skupina je jednoznačně určena těmito dílčími výběry.


\poznamka{Průběh: 0–9 min výběr proti funkcím; 9–19 min odvození a P1; 19–32 min P2–P3 a jeden z P4–P5; 32–42 min C1–C3; 42–45 min kontrola a shrnutí. C4–C5 a neprobraný vzorový příklad slouží k domácímu dokončení nebo jako zásoba. Smysl dělitele má přednost před počtem vypočtených úloh.}

\blok{Společně řešené příklady}

% Zdroj P1: import_hodiny/35_Kombinace_bez_opakovani.tex, Příklad 1.
\begin{prikladReseniBox}[title={Příklad P1}]
Kolika způsoby vybereme ze skupiny devíti různých osob tři členy komise bez rozlišených funkcí?
\tcblower
Pořadí tří vybraných osob nerozlišujeme. Uspořádaných výběrů by bylo $9\cdot8\cdot7$, ale každou komisi bychom tak započetli $3!=6$-krát. Proto
\[
\binom93=\frac{9\cdot8\cdot7}{3\cdot2\cdot1}=84.
\]
\end{prikladReseniBox}


% Zdroj P2: import_hodiny/35_Kombinace_bez_opakovani.tex, Příklad 2.
\begin{prikladReseniBox}[title={Příklad P2}]
Kolik úseček určuje šest bodů, nejsou-li žádné dva totožné?
\tcblower
Každé dva různé body jsou krajními body právě jedné úsečky; prohození konců ji nemění. Proto $\binom62=6\cdot5/2=15$. Různé dvojice mají různé krajní body, tedy různé úsečky. Není nutné předpokládat, že žádné tři body neleží na přímce: úsečky se mohou překrývat a stále být různé.
\end{prikladReseniBox}


% Zdroj P3: import_hodiny/35_Kombinace_bez_opakovani.tex, Příklad 3.
\begin{prikladReseniBox}[title={Příklad P3}]
Kolik čtyřčlenných skupin s přesně dvěma dívkami a dvěma chlapci lze vybrat ze šesti dívek a pěti chlapců?
\tcblower
Nejprve vybereme dvě dívky: $\binom62=15$ možností. Pro každou takovou dvojici můžeme vybrat dva chlapce $\binom52=10$ způsoby. Každá komise má právě jednu dívčí a jednu chlapeckou část, proto součinem dostáváme $15\cdot10=150$. Pořadí osob ani pořadí kategorií nevytváří další skupiny.
\end{prikladReseniBox}


% Zdroj P4: import_hodiny/35_Kombinace_bez_opakovani.tex, Příklad 4.
\begin{prikladReseniBox}[title={Příklad P4}]
Kolik úhlopříček má konvexní desetiúhelník?
\tcblower
Dvojice vrcholů určují $\binom{10}{2}=45$ úseček. Každá je buď stranou, nebo úhlopříčkou. Stran je deset, takže úhlopříček je $45-10=35$. Kontrola: z každého vrcholu vedou úhlopříčky do sedmi vrcholů; součet $10\cdot7$ počítá každou dvakrát, tedy opět $35$.
\end{prikladReseniBox}


% Zdroj P5: import_hodiny/35_Kombinace_bez_opakovani.tex, Příklad 5.
\begin{prikladReseniBox}[title={Příklad P5}]
Ze skupiny osmi lidí vybereme tři, přičemž právě jeden ze dvou konkrétních lidí A a B musí být vybrán.
\tcblower
Rozdělíme výběry na komise s A bez B a s B bez A. V každém případě doplníme dvě osoby z ostatních šesti: $\binom62=15$. Případy jsou disjunktní a jiné nejsou dovoleny, proto $15+15=2\binom62=30$.
\end{prikladReseniBox}


\newpage
\blok{Samostatná práce a zásoba k procvičení}

U každé úlohy napiš, co tvoří jednu skupinu a proč nezáleží na pořadí. U C2 výslovně využij geometrickou podmínku, u C3 nejprve rozděl výběr podle kategorií.

% Zdroj C1: import_hodiny/35_Kombinace_bez_opakovani.tex, Procvičování, položka 1.
\begin{zadaniReseniBox}[title={Úloha C1}]
Kolika způsoby vybereme dva z osmi různých žáků do dvojice bez rozlišených rolí?
\tcblower
Pořadí žáků nemění dvojici. $\binom82=8\cdot7/2=28$.
\end{zadaniReseniBox}

% Zdroj C2: import_hodiny/35_Kombinace_bez_opakovani.tex, Procvičování, položka 2.
\begin{zadaniReseniBox}[title={Úloha C2}]
Kolik trojúhelníků určuje sedm bodů, žádné tři neleží na přímce?
\tcblower
Každá trojice různých bodů určuje právě jeden trojúhelník, protože žádné tři neleží na přímce. Proto $\binom73=7\cdot6\cdot5/6=35$.
\end{zadaniReseniBox}

% Zdroj C3: import_hodiny/35_Kombinace_bez_opakovani.tex, Procvičování, položka 3.
\begin{zadaniReseniBox}[title={Úloha C3}]
Kolika způsoby ze šesti žen a čtyř mužů vybereme tříčlennou skupinu s přesně jedním mužem?
\tcblower
Vybereme jednoho muže čtyřmi způsoby a dvě ženy $\binom62=15$ způsoby. Výběry spojíme, tedy $\binom41\binom62=4\cdot15=60$.
\end{zadaniReseniBox}

% Zdroj C4: import_hodiny/35_Kombinace_bez_opakovani.tex, Procvičování, položka 4.
\begin{zadaniReseniBox}[title={Úloha C4}]
Kolik možností má výběr dvou karet z pěti různých karet, pokud pořadí nehraje roli?
\tcblower
Každá dvojice karet se započítá jednou, nikoli v obou pořadích: $\binom52=5\cdot4/2=10$.
\end{zadaniReseniBox}

% Zdroj C5: import_hodiny/35_Kombinace_bez_opakovani.tex, Procvičování, položka 5.
\begin{zadaniReseniBox}[title={Úloha C5}]
Kolika způsoby z devíti různých hráčů vybereme čtyři, musí-li být předem určený hráč P mezi vybranými?
\tcblower
Hráč P je pevně vybrán. Doplníme ještě tři z ostatních osmi hráčů: $\binom83=8\cdot7\cdot6/6=56$. Dalším faktorem za P nenásobíme.
\end{zadaniReseniBox}


\poznamka{Typické chyby: ponechání počtu variací bez dělení $k!$; přidání faktoru $2!$ za dívčí a chlapeckou část skupiny; započítání stran mezi úhlopříčky. U P5 „právě jeden“ vylučuje současný výběr A a B. V hodině 36 procvičíme samotná kombinační čísla a jednoduché rovnice s nimi.}
\end{document}
